\section{Maximum-score localization}
\label{l:maxscore:section}

We now select a cell by maximizing its posterior mass weighted by
depth. The selected posterior again satisfies a dyadic mass bound,
with a slack that accounts for conditioning on the selection event.
We first construct the selection and estimate its entropy. We then
prove a block-information estimate that charges one half of the
incoming slack. The iteration closes because the cell volume certifies
more information per level than the leading cost of naming the cell;
the depth entropy controls the remaining slack.

Put \(n=d-1\), let \(\tau\) be uniform probability on
\[
 Q=[-1/(2\sqrt n),1/(2\sqrt n)]^n,\qquad
 s(z)=(z,\sqrt{1-\|z\|^2}).
\]
Assign dyadic boundaries consistently to obtain a nested Borel
partition.  The proof of Lemma~\ref{ld:cube-geometry} applies on this
closure: a depth-\(i\) cell has side \(2^{-i}/\sqrt n\) and
\(\tau\)-mass \(2^{-ni}\), the graph map is \(2\)-Lipschitz, and
its first \(n\) coordinates recover \(z\).
Set
\[
 k=\ell=\lfloor n/8\rfloor,\qquad \beta=n/2,
\]
and take \(n\) sufficiently large that \(k=\ell\ge2\).
For a probability \(\mu\) supported on a depth-\(i\) cell \(C\),
the mass bound with slack \(a\ge0\) is
\begin{equation}\label{l:maxscore:spread}
 \mu(B)\le e^a2^{-\beta j}
 \quad\text{for every depth-\((i+j)\) subcell }B\subseteq C,
 \quad j\ge0.
\end{equation}

\subsection{Selecting a scale and paying for its label}

After a block, we restore this bound by examining all depths on the
cell chain of the sampled point. We reveal the first maximizing
depth and its cell. The resulting law is conditioned on the
selection event inside that cell, which is generally a proper subset.

\begin{lemma}[Maximum-score refinement]\label{l:maxscore:refinement}
Let a positive-probability discrete history, of probability \(P>0\),
have posterior \(\pi\) for \(Z\sim\tau\), supported on a depth-\(i\)
cell \(C\).  Thus \(\pi\le\tau/P\).  For \(z\in C\), let
\(C_{i+j}(z)\) be its depth-\((i+j)\) cell.  Choose the first
maximizer \(J\ge0\) of
\[
 2^{\beta j}\pi(C_{i+j}(Z)),\qquad j\ge0,
\]
and reveal \((J,B)\), where \(B=C_{i+J}(Z)\).
This is a measurable finite selection and
\[
 J\le\frac{\log(1/P)}{(n-\beta)\log2}.
\]
For an outcome \((j,B)\) of conditional probability \(p>0\),
the resulting posterior satisfies \eqref{l:maxscore:spread} in \(B\)
with slack \(a'=\log(\pi(B)/p)\ge0\).  Conditionally on the original
history,
\begin{align}
 \mathbb E a'&\le H(J)\le(1+\mathbb EJ)\log2,
 \label{l:maxscore:slack-charge}\\
 H(J,B)&\le(\beta+1)(\log2)\mathbb EJ+\log2.
 \label{l:maxscore:label-charge}
\end{align}
\end{lemma}

\begin{proof}
The score at depth zero is one.  Every score at depth \(j\) is at most
\[
 P^{-1}2^{-n(i+j)}2^{\beta j}
 =P^{-1}2^{-ni}2^{-(n-\beta)j},
\]
which tends to zero uniformly in \(z\).  Hence a maximum is attained
among finitely many depths.  The first such depth is measurable because
each score is constant on Borel cells and there are countably many
score comparisons.  Comparing its score with one proves the stated
bound on \(J\).

Let \(E_{j,B}\subseteq B\) be the actual selection event, of
\(\pi\)-mass \(p\).  Comparison with the depth-zero score gives
\(\pi(B)\ge2^{-\beta j}\).  If a further subcell
\(B'\subseteq B\) of relative depth \(m\) has positive posterior
mass after selection, it meets \(E_{j,B}\).  The maximizing inequality
at such a point implies
\(\pi(B')\le2^{-\beta m}\pi(B)\).  Its new posterior mass is
therefore at most \(p^{-1}\pi(B')\le e^{a'}2^{-\beta m}\);
if it has zero selected mass, the same bound holds.  Since
\(p\le\pi(B)\), the slack is nonnegative.

For fixed \(j\), let \(p_B\) be the probabilities of its cell
outcomes and \(p_j=\sum_Bp_B\).  These \(B\)'s are disjoint at
one grid level, so \(\sum_B\pi(B)\le1\).  The log-sum inequality
then gives
\[
 \sum_Bp_B\log\frac{\pi(B)}{p_B}\le p_j\log(1/p_j).
\]
Summing over \(j\) proves \(\mathbb E a'\le H(J)\).
Comparison of the law of \(J\) with the probability vector
\((2^{-j-1})_{j\ge0}\), by nonnegativity of relative entropy,
gives \(H(J)\le(1+\mathbb EJ)\log2\).  These quantities are
finite because \(J\) has the displayed finite bound.
Finally
\[
 \log(1/p)=\log(1/\pi(B))+a'\le\beta j\log2+a'.
\]
Averaging proves \eqref{l:maxscore:label-charge}.
\end{proof}

The refinement has two costs: the entropy of its label and the slack
in the selected posterior. The lemma bounds both in terms of the
expected depth increase. It remains to control the information in a
fresh block under \eqref{l:maxscore:spread}. The next estimate retains
the mass of a restricted event, so that it can subsequently be applied
to sets where a message creates a large posterior density.

\subsection{Projection moments of a restricted measure}

The inverse-distance calculation is related to the Gaussian projection
method of Sharan, Sidford, and Valiant \cite{SSV2019}. We give the
full estimate for exact block observations and arbitrary measurable
state rules. Restrictions are left unnormalized throughout.

\begin{lemma}[Restricted projection moment]\label{l:maxscore:projection}
Let \(\mu\) be a probability measure supported on a depth-\(i\)
cell \(C\subseteq Q\) satisfying \eqref{l:maxscore:spread}
with slack \(a\ge0\).
Set \(r=2^{-i}\).  There are absolute constants \(C_2,C_3\ge1\)
such that, with \(K=C_2^ne^a\), the following assertions hold.
For every \(\mu\)-measurable \(F\subseteq C\) of mass \(f=\mu(F)\)
and every affine subspace \(P\subseteq\mathbb R^d\) of dimension at
most \(\ell\),
\begin{equation}\label{l:maxscore:fractional}
 \int_F\operatorname{dist}(s(z),P)^{-k}\,d\mu(z)
 \le 2r^{-k}K^{1/3}f^{2/3}.
\end{equation}
Let \(\gamma_{k,d}\) be the law of a \(k\)-by-\(d\) standard
Gaussian matrix \(A\).  The joint image of
\(\gamma_{k,d}(dA)\,\mu|_F(dz)\) under
\((A,z)\mapsto(A,As(z))\) has a nonnegative density \(v_F(A,b)\)
relative to \(\gamma_{k,d}(dA)\,db\), and
\begin{equation}\label{l:maxscore:moment}
 \int v_F(A,b)^\ell\,d\gamma_{k,d}(A)\,db
 \le f\left[C_3^k r^{-k}K^{1/3}f^{2/3}\right]^{\ell-1}.
\end{equation}
In particular, for almost every \(A\), \(v_F(A,\cdot)\) is a density
of the image of the unnormalized measure \(\mu|_F\) under \(As\).
When \(f=0\), all these measures and integrals are zero.
\end{lemma}

\begin{proof}
Let \(z_c\) be the center of \(C\) and \(s_c=s(z_c)\).
The radius of \(C\) is \(r/2\), so its graph image lies in
\(B(s_c,r)\).  For \(0<h\le r\), choose \(j\ge0\) with
\(r2^{-j}\le h<2r2^{-j}\).  A depth-\((i+j)\) cell meeting
the first-coordinate projection of an ambient ball of radius \(h\)
lies in the concentric \(n\)-ball of radius \(3r2^{-j}\).
These cells have disjoint interiors and volume
\((r2^{-j}/\sqrt n)^n\).  The unit-ball volume estimate
\(v_n\le(2\pi e/n)^{n/2}\) therefore bounds their number by
\(C^n\), for an absolute \(C\), by the Gaussian-integral proof in
Lemma~\ref{ld:cube-geometry}.  By \eqref{l:maxscore:spread},
\begin{equation}\label{l:maxscore:ball}
 \mu\{z:\|s(z)-x\|\le h\}
 \le C^ne^a(h/r)^\beta.
\end{equation}
After increasing \(C\), this bound also holds trivially when \(h>r\).

If \(P\) has dimension \(m\le\ell\), the projection of
\(B(s_c,r)\) onto \(P\) is contained in an \(m\)-ball of radius
\(r\).  A maximal \(h\)-separated set in that ball has at most
\((1+2r/h)^m\) points: its disjoint radius-\(h/2\) balls fit
inside the concentric ball of radius \(r+h/2\).  Ambient balls of
radius \(2h\) about these centers cover points of \(B(s_c,r)\)
at distance at most \(h\) from \(P\).  Applying the ball bound and
absorbing the covering constants gives, for \(C_2\) absolute,
\begin{equation}\label{l:maxscore:tube}
 \mu\{z:\operatorname{dist}(s(z),P)\le h\}
 \le K(h/r)^{\beta-m}\qquad(0<h\le r).
\end{equation}
The estimate also shows that \(P\) has zero \(\mu\)-mass.

Suppose \(f>0\), and put \(\theta=k/(\beta-m)\).
Since \(m\le\ell=k\le n/8\), we have \(\theta\le1/3\).
For \(t\ge1\), the \(\mu\)-mass within \(F\) where
\((r/\operatorname{dist}(s(z),P))^k>t\) is at most
\(\min\{f,Kt^{-1/\theta}\}\).  Integrating this tail and splitting
at \(t=(K/f)^\theta\ge1\) gives
\[
 \int_F\left(\frac r{\operatorname{dist}(s(z),P)}\right)^k\,d\mu
 \le\frac{K^\theta f^{1-\theta}}{1-\theta}
 \le2K^{1/3}f^{2/3}.
\]
Here \(K\ge1\), \(0<f\le1\), and \(\theta\le1/3\) justify the
last inequality.  This proves \eqref{l:maxscore:fractional}.

Let \(\nu_A\) be the image of \(\mu|_F\) under \(z\mapsto As(z)\).
If \(F\) is measurable only in the completion, replace it by a Borel
representative modulo \(\mu\), which leaves the restricted measure
unchanged.
For \(h>0\), write \(v_k\) for the volume of the unit ball in
\(\mathbb R^k\) and define its ball-average density
\[
 v_{F,h}(A,b)=\frac{\nu_A(B(b,h))}{v_kh^k}.
\]
Expand the \(\ell\)th power using \(z_1,\ldots,z_\ell\) integrated
against \((\mu|_F)^{\otimes\ell}\), and put \(s_j=s(z_j)\).
The intersection of the balls \(B(As_j,h)\) has volume at most
\(v_kh^k\), and is empty unless
\(\|A(s_j-s_1)\|\le2h\) for all \(2\le j\le\ell\).
For a fixed tuple, condition successively on the preceding difference
images.  Orthogonal decomposition of each Gaussian row shows that the
conditional image at step \(j\) has an independent centered residual
with covariance
\(\operatorname{dist}(s_j,P_j)^2I_k\), where
\[
 P_j=s_1+\operatorname{span}\{s_t-s_1:2\le t<j\}.
\]
The subspace bound makes all zero-distance tuples null.  On every
remaining tuple, the conditional probability of a radius-\(2h\) ball
is at most
\[
 v_k(2h)^k(2\pi)^{-k/2}
 \operatorname{dist}(s_j,P_j)^{-k}.
\]
Multiplying these bounds, dividing by \((v_kh^k)^\ell\), and
including the intersection volume cancels every \(h\) and \(v_k\)
factor.  Integrating \(z_\ell,\ldots,z_2\) in this order and applying
\eqref{l:maxscore:fractional} at each step leaves the factor
\(\mu(F)=f\) for \(z_1\).  Absorbing the numerical Gaussian and
radius factors into \(C_3^k\) proves
\[
 \int v_{F,h}^\ell\,d\gamma_{k,d}\,db
 \le f[C_3^kr^{-k}K^{1/3}f^{2/3}]^{\ell-1}
\]
uniformly for \(h>0\).

For every continuous compactly supported test function on the matrix
and label space, integrals against these ball averages converge as
\(h\downarrow0\) to its integral against the exact joint image
measure.  This follows by averaging translations of the test function
over a shrinking ball and then using bounded convergence.  H\"older's
inequality bounds the limiting functional on
\(L^{\ell/(\ell-1)}(\gamma_{k,d}(dA)\,db)\) by the \(\ell\)th
root of the last display.  Continuous compactly supported tests are
dense for this locally finite reference measure.  Extension and
\(L^p\) duality therefore give a jointly measurable representing
density with \eqref{l:maxscore:moment}; equality on these tests
identifies it with the exact image measure and makes it nonnegative.
Finally, testing the joint identity on a countable algebra of rational
label boxes and arbitrary measurable matrix sets identifies the
conditional density for almost every \(A\).  The averaging operation
has only proved a statement about exact labels; it is absent from the
observation rule.
\end{proof}

\subsection{The one-half cost of slack}

\begin{lemma}[Information from one block]\label{l:maxscore:block}
In the setting of Lemma~\ref{l:maxscore:projection}, let \(Z\sim\mu\)
be independent of \(A\).  Suppose a variable \(U\), taking at most
\(W\le(k+2)2^{d^2}\) values, is generated from the exact data
\((A,As(Z))\) by a measurable probability kernel.  There is an
absolute \(C_0\) such that, for all sufficiently large \(d\),
\begin{equation}\label{l:maxscore:block-information}
 I(Z;U)\le C_0d+\frac a2.
\end{equation}
\end{lemma}

\begin{proof}
Write \(q_u(A,b)\) for the kernel probabilities, so that
\(q_u\ge0\) and \(\sum_uq_u=1\).  With the center \(s_c\)
used above, define
\[
 E=\{\|As(Z)-As_c\|\le3r\sqrt k\},\qquad
 V=v_k(3r\sqrt k)^k\le(3\sqrt{2\pi e}\,r)^k.
\]
For each fixed \(Z\), the vector \(A(s(Z)-s_c)\) is Gaussian
with covariance at most \(r^2I_k\).  Hence
\[
 \Pr(E^c\mid Z)\le e^{-9k/4}2^{k/2}\le e^{-k},
\]
by Markov's inequality and
\(\mathbb E\exp(\|G\|^2/4)=2^{k/2}\) for standard Gaussian
\(G\in\mathbb R^k\).
Define comparison probabilities
\[
 \lambda_u=\mathbb E_A\frac1V
       \int_{B(As_c,3r\sqrt k)}q_u(A,b)\,db,
 \qquad \sum_u\lambda_u=1.
\]
These probabilities use an auxiliary label uniform on the displayed
ball, independently conditional on \(A\).

For a measurable \(F\) of mass \(f\), apply H\"older to its exact
joint density \(v_F\) and the function
\(q_u(A,b)\mathbf1_{\{\|b-As_c\|\le3r\sqrt k\}}\).
Since \(0\le q_u\le1\), its conjugate norm is at most
\((V\lambda_u)^{1-1/\ell}\).  Lemma~\ref{l:maxscore:projection}
therefore gives
\[
 \Pr(Z\in F,U=u,E)
 \le (V\lambda_u)^{1-1/\ell}
       \{f[C_3^kr^{-k}K^{1/3}f^{2/3}]^{\ell-1}\}^{1/\ell}.
\]
The \(r\) powers cancel.  The exact exponent of \(f\) is
\[
 \frac{1+(2/3)(\ell-1)}{\ell}
 =\frac23+\frac1{3\ell}.
\]
Because \(f\le1\), it may be weakened to \(2/3\).
The exponent of \(K\) is \((\ell-1)/(3\ell)\le1/3\).
Thus, for an absolute \(C_4\ge1\),
\begin{equation}\label{l:maxscore:event}
 \Pr(Z\in F,U=u,E)
 \le L_0\lambda_u^{1-1/\ell}f^{2/3},
 \qquad L_0=C_4^ne^{a/3}.
\end{equation}

Put \(p_u=\Pr(U=u,E)\).  For \(p_u>0\), let
\(\nu_u=\mathcal L(Z\mid U=u,E)\) and \(f_u=d\nu_u/d\mu\).
Applying \eqref{l:maxscore:event} with \(F=C\) shows that
\[
 D_u^*=\frac{L_0\lambda_u^{1-1/\ell}}{p_u}\ge1;
\]
in particular \(\lambda_u>0\).  Since
\(\mu\{f_u>e^t\}\le e^{-t}\), a second application to that set
gives
\[
 \nu_u\{\log f_u>t\}\le\min\{1,D_u^*e^{-2t/3}\}
 \qquad(t\ge0).
\]
Integrating this upper bound gives
\(\tfrac32(\log D_u^*+1)\).  Relative entropy is no greater than
the expected positive part of the log density, so
\begin{equation}\label{l:maxscore:posterior-kl}
 D(\nu_u\Vert\mu)\le\frac32(\log D_u^*+1).
\end{equation}
This calculation is why the \(a/3\) in \(L_0\) becomes \(a/2\).

Let \(p_E=\sum_up_u\).  The log-sum inequality and the entropy
bound for at most \(W\) letters give
\[
 \sum_up_u\log\frac{\lambda_u}{p_u}\le p_E\log(1/p_E),
 \qquad
 \sum_up_u\log(1/p_u)\le p_E\log(1/p_E)+p_E\log W.
\]
The identity
\[
 \log\frac{\lambda_u^{1-1/\ell}}{p_u}
 =\left(1-\frac1\ell\right)\log\frac{\lambda_u}{p_u}
       +\frac1\ell\log(1/p_u)
\]
therefore bounds the weighted sum of these log ratios by
\(1+(\log W)/\ell\).  All zero terms are omitted.
For the outcomes \((U,E^c)\), their weighted divergences to \(\mu\)
sum to at most
\[
 p_{E^c}\log(1/p_{E^c})+p_{E^c}\log W.
\]
Indeed, conditioning first on \(E^c\) multiplies the density by
at most \(1/p_{E^c}\), and revealing \(U\) afterwards adds at
most its conditional entropy.  This bound is zero when \(p_{E^c}=0\).

The mutual information of the finite pair \((U,\mathbf1_E)\) is
the sum of these good and bad weighted divergences.  Combining the
last bounds with \eqref{l:maxscore:posterior-kl} gives
\[
 I(Z;U)\le I(Z;U,\mathbf1_E)
 \le Cn+\frac a2+C\frac{\log W}{\ell}
          +e^{-k}\log W+C.
\]
Here \((\log W)/\ell=O(d)\) and
\(e^{-k}\log W=O(1)\) in the stated width range.  Absorbing the
absolute terms proves \eqref{l:maxscore:block-information}.
\end{proof}

\subsection{Iteration and the accuracy bound}

Use the finite-state observation model of Section~\ref{l:model}, with
an integer horizon \(T\ge0\).

\begin{proposition}[Maximum-score information bound]
\label{l:maxscore:iteration}
Suppose \(M\le d^2\).  For \(N=\lceil T/k\rceil\), the learner
under \(Z\sim\tau\) has an augmented discrete history
\(\mathcal H_N\) satisfying
\begin{equation}\label{l:maxscore:total-information}
 I(Z;\mathcal H_N)\le C_6dN,\qquad
 I(Z;\widehat S)\le C_6dN
\end{equation}
for an absolute \(C_6\) and all sufficiently large \(d\).
If a data-independent shared seed selects the rules, the first
inequality holds conditional on that seed with the same constant;
the displayed output bound holds after averaging over it.
The statement holds for every finite \(T\).
\end{proposition}

\begin{proof}
First fix a realization of a data-independent shared seed, if present.
The graph derivative contains the identity in its first \(n\) rows, so
Lemma~\ref{l:prior-ac} applies to the cube prior.  For this fixed seed,
Lemma~\ref{l:borel} supplies everywhere-defined Borel rules with the
same prior law.  We do this separately for each seed value under
consideration; no jointly chosen family of Borel versions is needed.
Use the initial finite state as \(\mathcal H_0\).  By the joint
independence of the shared seed and initial state from the signal and
all rows in Section~\ref{l:model}, conditional on the fixed seed this
state has zero information about \(Z\), and the future rows remain
independent of both.  At an active block end record
the continuing state or the terminal state and its stopping position
within the block, including immediate stopping.  Given the preceding
history there are at most \((k+2)2^M\) such destinations.
After a stop, later messages may be dummy messages because the earlier
history already records the terminal state and index.  Padding the last
block by ignored Gaussian rows is harmless.  This count depends only on
the block length, not the total horizon.

Start at depth \(i_0=0\) with slack \(a_0=0\).
At step \(t\), apply Lemma~\ref{l:maxscore:block} to its outgoing
message \(U_t\).  Then apply Lemma~\ref{l:maxscore:refinement}
to the posterior given
\(\mathcal H_t^-=(\mathcal H_{t-1},U_t)\), and append its
selection \((J_t,B_t)\) to form \(\mathcal H_t\).
The new depth is \(i_t=i_{t-1}+J_t\), and its slack is \(a_t\).
The added labels use only \(Z\) and the history already observed.
Future Gaussian rows remain independent of that joint collection.
Given the starting state, the message is a kernel of the fresh block
data, so the block lemma applies on every positive-probability history.

Every such history of probability \(P\) has posterior at most
\(\tau/P\).  All entropies used below are finite.  Indeed a block
message has a finite alphabet, and the depth bound in
Lemma~\ref{l:maxscore:refinement} is integrable after averaging over
the preceding history because it is bounded by a constant times
\(\log(1/P)\).  The label-charge inequality then preserves finite
history entropy by induction.

The information chain rule, followed by
\eqref{l:maxscore:slack-charge} and \eqref{l:maxscore:label-charge},
gives
\[
 \begin{split}
 I(Z;\mathcal H_N)
 &\le\sum_{t=1}^N
 \left(C_0d+\tfrac12\mathbb E a_{t-1}
       +H(J_t,B_t\mid\mathcal H_t^-)\right)\\
 &\le C_5dN+
       (\beta+\tfrac32)(\log2)\sum_{t=1}^N\mathbb EJ_t
 \end{split}
\]
for an absolute \(C_5\).  On the other hand the final posterior
is supported on a depth-\(i_N\) cell, of prior mass \(2^{-ni_N}\).
Relative entropy to \(\tau\) is therefore at least
\(ni_N\log2\), since the conditional divergence to the uniform
law of that cell is nonnegative.  Averaging yields
\[
 I(Z;\mathcal H_N)\ge
 n(\log2)\mathbb Ei_N
 =n(\log2)\sum_{t=1}^N\mathbb EJ_t.
\]
The ratio \((\beta+3/2)/n=1/2+3/(2n)\) is bounded away from
one for large \(n\).  Substituting the last bound into the preceding
upper bound and moving that term to the left proves
\(I(Z;\mathcal H_N)\le C_6dN\).

The history records the terminal state and stopping index.  Conditional
on those data, the output uses only fresh randomness, so data processing
gives \(I(Z;\widehat S)\le I(Z;\mathcal H_N)\).
For each fixed seed, Lemma~\ref{l:borel} preserved the original
conditional prior law, so the output-information bound just proved
holds for that original law.  Integrating the original measurable
conditional information values over a shared seed \(R\), whose law
is independent of \(Z\), gives
\[
 I(Z;\widehat S)\le I(Z;\widehat S,R)
 =\mathbb E_R I(Z;\widehat S\mid R)\le C_6dN.
\]
When \(T=0\), the same argument has no blocks and the output is
independent of \(Z\).  This proves the stated zero upper bound as well.
\end{proof}

For a fixed unit output, angular success implies that \(Z\) lies
within \(\epsilon\) of its first \(n\) coordinates.  The ball-volume
estimate used above and the cube volume \(n^{-n/2}\) bound this
prior mass by \((\sqrt{2\pi e}\,\epsilon)^n\le(5\epsilon)^n\).
If the learner succeeds with probability at least \(2/3\) for every
unit signal, its cube success has the same lower bound.  Applying
\eqref{l:binary} to the success event under the joint and product
laws, and then using \eqref{l:maxscore:total-information}, gives
\begin{equation}\label{l:maxscore:accuracy}
 \frac23 n\log\frac1{5\epsilon}-\log2
 \le C_6d\lceil T/k\rceil.
\end{equation}
For \(L=\log(1/\epsilon)\ge\log10\),
\[
 L-\log5\ge
 \left(1-\frac{\log5}{\log10}\right)L.
\]
Thus \(L\le C(T/k+1)\) for large \(d\), uniformly in
\(0<\epsilon\le1/10\).  If \(L\ge2C\), then
\(T\ge k(L/C-1)\ge kL/(2C)\ge cdL\).

For the remaining values \(L<2C\), the original learner needs
\(T>d/10\) samples. Average its every-signal guarantee over a
uniform signal \(S\) on the full sphere, and grant it all \(T\)
pregenerated rows, including unused rows. If \(T\le d/10\),
Lemma~\ref{l:caps-endpoint} applies to this original randomized
learner. Writing \(V\) for the row span and letting \(r_0\)
decrease to \(\epsilon\) in \eqref{l:raw-success} gives
\begin{equation}\label{l:maxscore:reflection}
 \Pr(\text{success})
 \le\frac12+\frac12\Pr\{\|P_{V^\perp}S\|\le\epsilon\}
 \le\frac12+\frac{T}{2d(1-\epsilon^2)}.
\end{equation}
For \(T\le d/10\) and \(\epsilon\le1/10\), the right side is
less than \(2/3\).  Combining this linear bound with the large-\(L\)
bound proves \(T=\Omega(d\log(1/\epsilon))\).
The constant is absolute, and \(M(d)=o(d^2)\) ensures \(M\le d^2\)
eventually.  No restriction on \(T\) or on the rate of growth of \(L\)
was used.
